% Ch4: Sobolev Functions
% Evans & Gariepy, Chapter 4 (pp. 131--165)
% Lectures 10--12, Fall 2024

\section{October 7, 2024}

Rain all morning. The walk to class was miserable but at least the topic is good. We're starting Sobolev spaces, which is where analysis starts feeling like it has teeth.

\subsection{Weak Derivatives}

\begin{definition}[Weak derivative]
\label{def:weak-deriv}
Let $U \subseteq \R^n$ be open. Suppose $u, v \in L^1_{\mathrm{loc}}(U)$. We say $v$ is the \textbf{$\alpha$-th weak partial derivative} of $u$, written $D^\alpha u = v$, provided
\[
    \int_U u \, D^\alpha \varphi \, dx = (-1)^{|\alpha|} \int_U v \, \varphi \, dx
\]
for all test functions $\varphi \in C^\infty_c(U)$.
\end{definition}

Here's the idea. You move all the derivatives onto the smooth test function via integration by parts, then \emph{define} the derivative of $u$ to be whatever $v$ makes the identity work. If $u$ happens to be classically differentiable, you get the usual derivative back. But now functions with corners, kinks, jumps in their derivatives --- all fair game.

Weak derivatives are unique up to a.e.\ equivalence. That's the fundamental lemma of the calculus of variations, nothing more.


\subsection{Sobolev Spaces}

\begin{definition}[Sobolev space $W^{k,p}(U)$]
\label{def:sobolev-space}
For $k \in \mathbb{N}$ and $1 \leq p \leq \infty$, the Sobolev space $W^{k,p}(U)$ consists of all $u \in L^p(U)$ such that for each multi-index $\alpha$ with $|\alpha| \leq k$, the weak derivative $D^\alpha u$ exists and belongs to $L^p(U)$.
\end{definition}

The norm:
\[
    \|u\|_{W^{k,p}(U)} =
    \begin{cases}
        \displaystyle\left( \sum_{|\alpha| \leq k} \int_U |D^\alpha u|^p \, dx \right)^{1/p} & 1 \leq p < \infty, \\[8pt]
        \displaystyle\sum_{|\alpha| \leq k} \operatorname*{ess\,sup}_U |D^\alpha u| & p = \infty.
    \end{cases}
\]

$W^{k,p}(U)$ is a Banach space. For $p = 2$ it's a Hilbert space, usually written $H^k(U)$.

$W^{k,p}_0(U)$ is the closure of $C^\infty_c(U)$ in $W^{k,p}(U)$. Think of these as Sobolev functions that ``vanish on $\partial U$'' in some appropriate sense.

Product rules and chain rules work. If $u \in W^{1,p}(U)$ and $\Phi \in C^1$ with $\Phi'$ bounded, then $\Phi \circ u \in W^{1,p}(U)$ and $D(\Phi \circ u) = \Phi'(u) Du$.

One thing that confused me early on: $W^{k,p}(U)$ is \emph{not} just ``$C^k$ functions completed in the $W^{k,p}$ norm.'' That completion starting from $C^\infty_c$ gives you $W^{k,p}_0$, which is a different space. You genuinely need the weak derivative definition.


\subsection{Approximation by Smooth Functions}

The mollification proof is the template for everything in this chapter. Learn it once and you're set.

\begin{definition}[Mollifier]
Let $\eta \in C^\infty_c(\R^n)$ be the standard mollifier with $\eta \geq 0$, $\operatorname{spt} \eta \subseteq B(0,1)$, and $\int \eta = 1$. Set $\eta_\varepsilon(x) = \varepsilon^{-n} \eta(x/\varepsilon)$.
\end{definition}

For $u \in L^1_{\mathrm{loc}}(U)$, the mollification is $u^\varepsilon = \eta_\varepsilon * u$ on $U_\varepsilon = \{x \in U : \operatorname{dist}(x, \partial U) > \varepsilon\}$.

\begin{theorem}[Local approximation]
\label{thm:local-approx}
If $u \in W^{k,p}(U)$ with $1 \leq p < \infty$, then $u^\varepsilon \in C^\infty(U_\varepsilon)$ and $u^\varepsilon \to u$ in $W^{k,p}_{\mathrm{loc}}(U)$ as $\varepsilon \to 0$.
\end{theorem}

\begin{proof}[Proof sketch]
The trick is showing $D^\alpha(\eta_\varepsilon * u) = \eta_\varepsilon * D^\alpha u$ for $|\alpha| \leq k$ on $U_\varepsilon$. Write out the convolution, differentiate under the integral (legal since $\eta_\varepsilon$ is smooth and compactly supported), recognize the result as $(-1)^{|\alpha|} \int u \, D^\alpha_y \eta_\varepsilon(\cdot - y) \, dy$, then use the definition of weak derivative to flip the derivatives back. Once you have this identity, $L^p$ convergence is the standard mollification result.
\end{proof}

Now the local-to-global step.

\begin{theorem}[Meyers--Serrin]
\label{thm:meyers-serrin}
$C^\infty(U) \cap W^{k,p}(U)$ is dense in $W^{k,p}(U)$ for $1 \leq p < \infty$.
\end{theorem}

The proof uses a partition of unity subordinate to an exhaustion of $U$. You mollify each piece with a small enough $\varepsilon$ and sum. It's a careful $\varepsilon/2^k$ argument, nothing deep, but you have to track the supports.

Can you approximate by $C^\infty(\overline{U})$ functions? In general, no.

\begin{theorem}[Approximation up to the boundary]
\label{thm:global-approx-boundary}
If $\partial U$ is Lipschitz, then $C^\infty(\overline{U})$ is dense in $W^{k,p}(U)$ for $1 \leq p < \infty$.
\end{theorem}

The Lipschitz condition lets you locally flatten the boundary and push the function slightly inward, where you can mollify safely. Without it, you can build domains where the approximation fails.

Three approximation theorems, each stronger, each needing more hypotheses. I had to draw a diagram to keep them straight.



\section{October 9, 2024}

Office hours ran long yesterday so I didn't get to review traces before today. Walking in blind.

\subsection{Traces}

For smooth functions $u \in C^\infty(\overline{U})$, restricting to $\partial U$ is trivial. For Sobolev functions you can't just ``evaluate on the boundary'' because they're only defined a.e. The trace operator fixes this.

\begin{theorem}[Trace theorem]
\label{thm:trace}
Assume $U$ is bounded and $\partial U$ is Lipschitz. There exists a bounded linear operator
\[
    T : W^{1,p}(U) \to L^p(\partial U)
\]
such that $Tu = u|_{\partial U}$ for $u \in W^{1,p}(U) \cap C(\overline{U})$, and $\|Tu\|_{L^p(\partial U)} \leq C \|u\|_{W^{1,p}(U)}$.
\end{theorem}

The trace characterizes $W^{1,p}_0(U)$: a function $u \in W^{1,p}(U)$ belongs to $W^{1,p}_0(U)$ if and only if $Tu = 0$ on $\partial U$. So $W^{1,p}_0$ really is the space of Sobolev functions vanishing at the boundary, exactly as the notation suggests.

The sharper statement is that $T$ maps into the fractional Sobolev space $W^{1-1/p, p}(\partial U)$, but Evans--Gariepy don't go there.


\subsection{Extensions}

\begin{theorem}[Extension]
\label{thm:extension}
Assume $U$ is bounded and $\partial U$ is Lipschitz. For each $1 \leq p \leq \infty$, there exists a bounded linear operator
\[
    E : W^{1,p}(U) \to W^{1,p}(\R^n)
\]
such that $Eu = u$ a.e.\ in $U$, $\operatorname{spt}(Eu)$ is bounded, and
\[
    \|Eu\|_{W^{1,p}(\R^n)} \leq C \|u\|_{W^{1,p}(U)}.
\]
\end{theorem}

Near the boundary, reflect the function across $\partial U$. Because $\partial U$ is Lipschitz, you can locally flatten it, reflect, then patch with a partition of unity. Multiply by a cutoff for bounded support.

Extensions matter because lots of results are easier to prove on all of $\R^n$ with no boundary issues. Then you transfer back via $E$. The Sobolev inequalities are the prime example.


\subsection{Gagliardo--Nirenberg--Sobolev ($p < n$)}
\label{thm:GNS}

This is the heart of the chapter. The notation is brutal and I won't apologize for it.

The exponent $p^* = np/(n-p)$ looks arbitrary until you do dimensional analysis. $\|u\|_{L^q}$ has dimensions of $(\text{length})^{n/q}$ times pointwise size, and $\|Du\|_{L^p}$ has dimensions of $(\text{length})^{n/p - 1}$ times pointwise size. For the inequality to be scale-invariant, we need $n/q = n/p - 1$, giving $q = np/(n-p) = p^*$. That's the only exponent that works.

\begin{theorem}[Gagliardo--Nirenberg--Sobolev]
\label{thm:gns}
Assume $1 \leq p < n$. There exists $C = C(p,n)$ such that
\[
    \|u\|_{L^{p^*}(\R^n)} \leq C \|Du\|_{L^p(\R^n)}
\]
for all $u \in C^1_c(\R^n)$, where $p^* = \dfrac{np}{n-p}$.
\end{theorem}

The right-hand side only involves $Du$, not $u$ itself. That's because we're on all of $\R^n$ with compactly supported functions.

By density this extends to $W^{1,p}(\R^n)$, and via the extension operator to bounded domains with Lipschitz boundary:
\[
    \|u\|_{L^{p^*}(U)} \leq C \|u\|_{W^{1,p}(U)}.
\]

So $W^{1,p}(U) \hookrightarrow L^{p^*}(U)$ continuously. You trade one derivative for better integrability.

\begin{corollary}[Poincar\'e inequality]
\label{cor:poincare}
If $U$ is bounded and $1 \leq p < n$, then for $u \in W^{1,p}_0(U)$:
\[
    \|u\|_{L^q(U)} \leq C \|Du\|_{L^p(U)} \quad \text{for all } 1 \leq q \leq p^*.
\]
\end{corollary}


\subsection{Morrey's Inequality ($p > n$)}

Morrey's inequality is what makes $p > n$ special. You get continuity for free.

\begin{theorem}[Morrey's inequality]
\label{thm:morrey}
Assume $n < p \leq \infty$. There exists $C = C(p,n)$ such that
\[
    \|u\|_{C^{0,\gamma}(\R^n)} \leq C \|Du\|_{L^p(\R^n)}
\]
for all $u \in C^1_c(\R^n)$, where $\gamma = 1 - n/p$.
\end{theorem}

So $W^{1,p}(U) \hookrightarrow C^{0,1-n/p}(\overline{U})$ for $p > n$. One derivative in $L^p$ with $p$ large enough, and the function \emph{has} to be H\"older continuous. As $p \to \infty$, $\gamma \to 1$ and you approach Lipschitz. As $p \to n^+$, $\gamma \to 0$ and you just barely get continuity.

The proof integrates $|u(x) - u(y)|$ along line segments, switches to polar coordinates, and applies H\"older's inequality. The key geometric observation: the average of $|Du|$ over a ball controls the oscillation of $u$.


\subsection{The borderline $p = n$}

The case $p = n$ is subtle. You don't get $L^\infty$ (that would be Morrey, but $p = n$ doesn't satisfy $p > n$), and $p^* = np/(n-p)$ blows up. What you get is $W^{1,n}(U) \hookrightarrow L^q(U)$ for all $q < \infty$, but \emph{not} $q = \infty$. Almost there, but not quite.


\subsection{Embedding summary}

For $U$ bounded with Lipschitz boundary:
\begin{center}
\renewcommand{\arraystretch}{1.3}
\begin{tabular}{lll}
    \toprule
    \textbf{Regime} & \textbf{Embedding} & \textbf{Key feature} \\
    \midrule
    $1 \leq p < n$ & $W^{1,p}(U) \hookrightarrow L^{p^*}(U)$ & Improved integrability \\
    $p = n$ & $W^{1,n}(U) \hookrightarrow L^q(U)$, all $q < \infty$ & Almost $L^\infty$ \\
    $p > n$ & $W^{1,p}(U) \hookrightarrow C^{0,1-n/p}(\overline{U})$ & H\"older continuity \\
    \bottomrule
\end{tabular}
\end{center}

The dimension $n$ is the threshold. Everything turns on whether $p$ is below, at, or above $n$.



\section{October 14, 2024}

I think I finally understand the Sobolev embeddings. It only took three passes through the proofs. My handwritten notes from the first pass are unreadable garbage. Today we do compactness, capacity, and the fine properties of Sobolev functions.

\subsection{Rellich--Kondrachov Compactness}
\label{thm:rellich-kondrachov}

This is Arzel\`a--Ascoli for Sobolev spaces.

\begin{theorem}[Rellich--Kondrachov]
\label{thm:rellich}
Assume $U \subset \R^n$ is bounded and open with Lipschitz boundary. Then:
\begin{enumerate}
    \item If $1 \leq p < n$: $W^{1,p}(U) \hookrightarrow\hookrightarrow L^q(U)$ for each $1 \leq q < p^*$.
    \item If $p = n$: $W^{1,p}(U) \hookrightarrow\hookrightarrow L^q(U)$ for each $1 \leq q < \infty$.
    \item If $p > n$: $W^{1,p}(U) \hookrightarrow\hookrightarrow C^{0,\beta}(\overline{U})$ for each $0 < \beta < 1 - n/p$.
\end{enumerate}
Here $\hookrightarrow\hookrightarrow$ means compact embedding: bounded sequences have convergent subsequences.
\end{theorem}

Notice the strict inequalities. You get compactness into $L^q$ for $q < p^*$, but \emph{not} at the critical exponent $q = p^*$. The embedding $W^{1,p} \hookrightarrow L^{p^*}$ is continuous but not compact. This matters enormously in the calculus of variations --- if you've ever heard someone say ``loss of compactness at the critical exponent,'' this is what they mean.

\begin{proof}[Proof idea for case (1)]
Take a bounded sequence $(u_k)$ in $W^{1,p}(U)$. Extend to $\R^n$, then mollify. For fixed $\varepsilon > 0$, the mollified sequence $(u_k^\varepsilon)$ is bounded in $C^1$ on compact sets (from the $W^{1,p}$ bound and properties of convolution), so Arzel\`a--Ascoli gives a convergent subsequence. Diagonal argument over $\varepsilon \to 0$. The $q < p^*$ restriction enters because you need spare integrability to interpolate between $L^1$ convergence of mollifications and the $L^{p^*}$ bound.
\end{proof}

I won't pretend I understand every detail of the interpolation step. The mechanism is: Sobolev regularity gives equicontinuity after mollification, Arzel\`a--Ascoli gives convergence, interpolation gives $L^q$ convergence for subcritical $q$.


\subsection{Capacity}

Capacity confused me for weeks. Here's what I eventually understood: sets of capacity zero are the sets that Sobolev functions can't ``see.'' Measure zero sets are invisible to $L^p$ functions. Capacity zero sets are invisible to Sobolev functions. But capacity zero is a \emph{finer} notion --- there are measure-zero sets with positive capacity.

\begin{definition}[$p$-capacity]
\label{def:p-capacity}
For $1 \leq p < n$ and $A \subseteq \R^n$,
\[
    \operatorname{cap}_p(A) = \inf\left\{ \int_{\R^n} |Du|^p \, dx : u \in W^{1,p}(\R^n), \; u \geq 1 \text{ on a neighborhood of } A \right\}.
\]
\end{definition}

Properties:
\begin{enumerate}
    \item Monotone: $A \subseteq B \implies \operatorname{cap}_p(A) \leq \operatorname{cap}_p(B)$.
    \item Countably subadditive: $\operatorname{cap}_p\bigl(\bigcup_k A_k\bigr) \leq \sum_k \operatorname{cap}_p(A_k)$.
    \item If $\operatorname{cap}_p(A) = 0$, then $\mathcal{H}^s(A) = 0$ for all $s > n - p$.
    \item $\operatorname{cap}_p(A) = 0 \implies \mathcal{L}^n(A) = 0$, but not conversely.
\end{enumerate}

The relationship to Hausdorff measure is the one that matters most in practice. It tells you the dimension of the exceptional set where Sobolev functions can misbehave.


\subsection{Quasicontinuity and Precise Representatives}

A Sobolev function is only defined a.e., so talking about pointwise values is dangerous. But you can choose a representative that's well-behaved outside an arbitrarily small set in the capacity sense.

\begin{definition}[Quasicontinuous]
\label{def:quasicontinuous}
A function $u : U \to \R$ is \textbf{$p$-quasicontinuous} if for every $\varepsilon > 0$ there exists an open set $V$ with $\operatorname{cap}_p(V) < \varepsilon$ such that $u|_{U \setminus V}$ is continuous.
\end{definition}

\begin{theorem}[Quasicontinuous representative]
\label{thm:quasi-rep}
Every $u \in W^{1,p}(\R^n)$ ($1 \leq p \leq n$) has a $p$-quasicontinuous representative $\tilde{u}$, unique up to sets of $\operatorname{cap}_p$-zero.
\end{theorem}

Sobolev functions are ``almost continuous.'' The discontinuity set can be made as small as you want in the capacity sense. Not continuous. Not just measurable. Quasicontinuous --- that's the right notion of regularity for these functions.

The \textbf{precise representative} is defined pointwise:
\[
    u^*(x) = \lim_{r \to 0} \fint_{B(x,r)} u \, dy,
\]
whenever this limit exists. For $u \in W^{1,p}$, $u^*$ exists and equals $\tilde{u}$ outside a set of $\operatorname{cap}_p$-zero.

This connects back to the differentiation of measures stuff from the Vitali covering lectures. Same idea of recovering pointwise information from integral averages, but now the exceptional set is measured by capacity instead of Hausdorff measure.


\subsection{Differentiability on Lines}

Here's an equivalent way to think about Sobolev functions that avoids integration by parts entirely.

\begin{theorem}
\label{thm:diff-lines}
If $u \in W^{1,p}(U)$ for $1 \leq p \leq \infty$, then for each $i = 1, \ldots, n$, the function $t \mapsto u(x_1, \ldots, x_{i-1}, t, x_{i+1}, \ldots, x_n)$ is absolutely continuous on the relevant line segment for $\mathcal{L}^{n-1}$-a.e.\ choice of the remaining variables. Its classical derivative equals the weak derivative $D_i u$ a.e.
\end{theorem}

The converse is also true: if $u \in L^p(U)$ is absolutely continuous on a.e.\ line in each coordinate direction, and the classical partials are in $L^p$, then $u \in W^{1,p}(U)$.

Sobolev functions are exactly the $L^p$ functions that are absolutely continuous on almost every axis-parallel line. Clean characterization. Brian asked me why you can't just use this as the definition from the start. I think you could, honestly, but the integration-by-parts version plays better with the PDE applications.
